\documentclass[a4paper,oneside,12pt]{amsart}

%\usepackage[spanish, activeacute]{babel}
\usepackage[utf8]{inputenc}
%\usepackage[latin1]{inputenc}
\usepackage{latexsym}
\usepackage{multicol}
\usepackage{enumerate}
\usepackage{enumitem}
\usepackage[heightrounded]{geometry}
\geometry{top=2cm,bottom=2cm,left=2cm,right=2cm}
\usepackage{graphicx}

\DeclareMathOperator{\cond}{cond}
\DeclareMathOperator{\Id}{Id}
\DeclareMathOperator{\re}{Re}
%\DeclareMathOperator{\im}{Im}
\newcommand{\I}{\mathbb{I}}
\newcommand{\R}{\mathbb{R}}
\newcommand{\Q}{\mathbb{Q}}
\newcommand{\C}{\mathbb{C}}
\newcommand{\Z}{\mathbb{Z}}
\newcommand{\N}{\mathbb{N}}
\newcommand{\xopt}{x^{*}}
\newcommand{\yopt}{y^{*}}
\DeclareMathOperator{\im}{Im}

\newcommand{\nombremateria}
{Investigaci\'on Operativa}
\newcommand{\nombrecuatrimestre}
    {Segundo cuatrimestre de 2025}
\newcommand{\numeroparcial}
    { Recuperatorio del Segundo Parcial }
\newcommand{\fecha}
    {16 de Diciembre de 2025}
\newcommand{\numerotema}
    {\bf 1}

%\oddsidemargin -1cm \evensidemargin -1cm \textwidth 17.5cm
%topmargin 1.2cm \textheight 25cm
%\setlength{\textheight}{25cm}

\def \ds{\displaystyle}
\theoremstyle{definition}
\newtheorem{quest}{Ejercicio}
\thispagestyle{empty}

\begin{document}

\hfill \hfill
\begin{tabular}[c]{|c|c|c|c|}
\hline 
1 & 2 & 3 & Calificaci\'on \\ \hline \hline
\hspace{1.2 cm} & \hspace{1.2 cm} & \hspace{1.2 cm} &\hspace{1.5 cm} \\
\hspace{1.2 cm} & \hspace{1.2 cm} & \hspace{1.2 cm} &\hspace{1.5 cm} \\ \hline

\end{tabular}

\vspace{-1cm}

\noindent Nombre y apellido:

%\vspace{0.2cm}

\noindent Número de libreta:

\noindent \hrulefill 

%\smallskip
\renewcommand{\baselinestretch}{1.1}

\begin{center}
	\large \textbf{\nombremateria} 
	
	\numeroparcial -- \fecha
\end{center}

%\vspace{.6cm}
\renewcommand{\theenumi}{\textbf{\arabic{enumi}}}

\begin{center}
    \textbf{Entregar todos los ejercicios en hojas separadas.}
\end{center}
%%% EMPIEZA EL EJERCICIO 1 %%%

\begin{quest} 
\verb|[5.5 pts.]| Dado el siguiente problema de Programación Lineal:
    \[
        \begin{array}{rrcrcrcr}
            \max & 2x_1 & + & x_2 & + & x_3 & & \\
            \text{s.a:} & x_1 & + & x_2 & + & x_3 & \geq & 4\\
            & x_1 &  &    & + & 2x_3 & \leq & 6 \\
            & x_1 & - & x_2 & - & x_3 & = & -2\\
            & \multicolumn{5}{r}{x_1,x_2,x_3} & \geq & 0 \\
        \end{array}
    \]
    \begin{enumerate}[label=\alph*)]
        \item \verb|[3 pts.]| Resolver el problema utilizando el Método Big-M.
        \item \verb|[1 pt.]| Plantear el problema dual.
        \item \verb|[1.5 pts.]| A partir de la solución óptima obtenida en el ítem a), hallar una solución
        óptima del problema dual.
    \end{enumerate}

\end{quest}

%%% TERMINA EL EJERCICIO 1 %%%

%%% EMPIEZA EL EJERCICIO 2 %%%

\begin{quest}
    \verb|[1 pt.]| Los siguientes diccionarios corresponden a problemas de Programación Lineal con
    objetivo de maximizar. Utilizando el criterio usual de elección de qué variable entra a la
    base, determinar y justificar qué podemos deducir a partir de cada uno de ellos:
    \begin{multicols}{2}
        \begin{enumerate}[label=\alph*)]
            \item Única solución óptima no degenerada.
            \item Única solución óptima degenerada.
            \item Solución degenerada.
            \item No acotado.
            \item Infinitas soluciones óptimas.
        \end{enumerate}
    \end{multicols}
    \begin{multicols}{2}
        \begin{itemize}
            \item[1)] \verb|[0.25 pts.]|
            \[
                \begin{array}{rrrrrrr}
                    x_2 & = & 2           & - & \frac{1}{2}x_1 & - & 2x_5            \\
                    x_3 & = &             & - & x_1            & + & x_5            \\
                    x_4 & = & \frac{1}{3} & + & x_1            &   &                \\ \hline
                    z   & = & -3          & + & x_1               & - & x_5
                \end{array}
            \]
            \item[3)] \verb|[0.25 pts.]|
            \[
                \begin{array}{rrrrrrr}
                    x_1 & = & 1 &   &                & + & \frac{1}{5}x_5 \\
                    x_4 & = & 1 & + & x_2            & - & x_5            \\
                    x_3 & = & 2 & - & \frac{1}{4}x_2 &   &                \\ \hline
                    z   & = &   &   &                & - & 3x_5
                \end{array}
            \]
            \item[2)] \verb|[0.25 pts.]|
            \[
                \begin{array}{rrrrrrr}
                    x_3 & = & 1           &  - & x_2                &   &               \\
                    x_1 & = & \frac{2}{3} & + & \frac{1}{2}x_2 & + & \frac{1}{2}x_5 \\
                    x_4 & = & 2           & - & 2x_2           & + & 4x_5           \\ \hline
                    z   & = & -4          & + & x_2 & + & 2x_5
                \end{array}
            \]
            \item[4)] \verb|[0.25 pts.]|
            \[
                \begin{array}{rrrrrrr}
                    x_4 & = & 1           & + & 2x_1 & + & \frac{1}{2}x_5  \\
                    x_2 & = & 5           & + & x_1  & - & 3x_5 \\
                    x_3 & = & \frac{2}{3} & - & x_1  & - & x_5  \\ \hline
                    z   & = & 1           & - & x_1  & - & x_5
                \end{array}
            \]

        \end{itemize}
    \end{multicols}

\end{quest}
%%% TERMINA EL EJERCICIO 2 %%%

%%% EMPIEZA EL EJERCICIO 3 %%%

\begin{quest} 
\verb|[3.5 pts.]| Utilizar el algoritmo de Branch \& Bound para resolver el siguiente problema de Programación Lineal
Entera, detallando cada división en subproblemas mediante el esquema de árbol visto en clase.
\[\begin{array}{rrrrrrrrc}
\max & 3x_1 & + & 4x_2 &  &   \\
\text{s.a:}	 & 2x_1 & + & x_2 & \leq & 14  \\
			 & 2x_1 & + & 3x_2 & \leq & 17 \\
			 &2x_1 & +  & 2x_2 & \leq & 15 \\
			 &	\multicolumn{3}{r}{x_1,x_2} & \geq & 0 \\
%			 & & & x_2 & \geq & 0 \\
			 & \multicolumn{3}{r}{x_1,x_2} & \in & \mathbb{Z}
\end{array}
\]
\end{quest}
%%% TERMINA EL EJERCICIO 3 %%%

%%% TERMINA EL PARCIAL %%%

\newpage

   \begin{tabular}{|l|l|} \hline
        \multicolumn{1}{|c|}{Primal} & \multicolumn{1}{c|}{Dual}  \\ \hline
        Obj: Minimizar & Obj: Maximizar \\ \hline
        % Obj: Minimizar & Obj: Maximizar \\ \hline
        $\imath$-ésima restricción $\leq$ & $\imath$-ésima variable $\leq 0$ \\
        $\imath$-ésima restricción $\geq$ & $\imath$-ésima variable $\geq 0$ \\
        $\imath$-ésima restricción $=$ & $\imath$-ésima variable libre \\
        $\jmath$-ésima variable $\geq 0$ & $\jmath$-ésima restricción $\leq$ \\
        $\jmath$-ésima variable $\leq 0$ & $\jmath$-ésima restricción $\geq$ \\
        $\jmath$-ésima variable libre & $\jmath$-ésima restricción $=$ \\ \hline
    \end{tabular}
    \begin{tabular}{|l|l|} \hline
        \multicolumn{1}{|c|}{Primal} & \multicolumn{1}{c|}{Dual}  \\ \hline
        Obj: Maximizar & Obj: Minimizar \\ \hline
        % Obj: Minimizar & Obj: Maximizar \\ \hline
        $\imath$-ésima restricción $\leq$ & $\imath$-ésima variable $\geq 0$ \\
        $\imath$-ésima restricción $\geq$ & $\imath$-ésima variable $\leq 0$ \\
        $\imath$-ésima restricción $=$ & $\imath$-ésima variable libre \\
        $\jmath$-ésima variable $\geq 0$ & $\jmath$-ésima restricción $\geq$ \\
        $\jmath$-ésima variable $\leq 0$ & $\jmath$-ésima restricción $\leq$ \\
        $\jmath$-ésima variable libre & $\jmath$-ésima restricción $=$ \\ \hline
    \end{tabular}
    \vspace{1cm}

    \textbf{Teorema débil de dualidad:} Para un problema con objetivo de maximizar, sean $(x_1,\cdots,x_n)$ solución
factible del primal e $(y_1,\cdots, y_m)$ solución factible del dual, entonces:
    \[ \displaystyle \sum_{j=1}^n c_jx_j \leq \sum_{i=1}^m b_iy_i\]
    (si el objetivo del primal es minimizar, se invierte la desigualdad).

    \textbf{Teorema fundamental de dualidad}: Si el problema primal tiene solución óptima $(x_1^{\ast},\cdots,x_n^{\ast})$, entonces el dual
    tiene solución óptima $(y_1^{\ast},\cdots,y_m^{\ast})$ tal que:
    \[ \displaystyle \sum_{j=1}^n c_jx_j^{\ast} = \sum_{i=1}^m b_iy_i^{\ast}\]

    \textbf{Teorema de Holgura Complementaria:} Sean $x^{*}=(x_1^{\ast},\cdots,x_n^{\ast})$ solución del primal e
        $y^{*}=(y_1^{\ast},\cdots,y_m^{\ast})$ solución del dual, las siguientes son condiciones necesarias y
suficientes para la optimalidad simultánea de $\xopt$ e $\yopt$:
    \[
    \begin{array}{ccccc}
        \displaystyle \sum_{i=1}^m a_{ij}\yopt_i=c_j & \text{o} & \xopt_j = 0 & \; \text{(o ambos)} \quad & \forall j=1,\cdots,n  \\
        \displaystyle \sum_{j=1}^m a_{ij}\xopt_j=b_i & \text{o} & \yopt_i = 0 & \; \text{(o ambos)} \quad & \forall i=1,\cdots,m  \\
    \end{array} \]
\end{document}

